\begin{textsample}{POS Dim 2 – es – Score 35.00 – es/es\_inf\_intro\_4.txt}  \label{ex:f2_pos_241}
EDUCAÇÃO BÁSICA E SUAS BASES LEGAIS Além do estudo profundo da Base Nacional Comum Curricular, a \textbf{equipe} de \textbf{currículo} realizou estudos dos documentos normativos e legais da educação nacional ( Constituição Federal de 1988, LDB 9394/96, Diretrizes Nacionais da Educação Básica : Diversidade e Inclusão de 2013 ), de \textbf{currículos} nacionais e internacionais e, principalmente, dos \textbf{currículos} já construídos e vividos na rede estadual, no caso o \textbf{Currículo} Básico Escola Estadual - CBEE ( ESPÍRITO SANTO, 2009 ), e nas redes \textbf{municipais} do Espírito Santo¹. ¹ Foram considerados os documentos curriculares enviados pelos \textbf{municípios} que compartilharam seus documentos a título de contribuição para construção do \textbf{Currículo} do Espírito Santo, sendo eles : Aracruz, Boa Esperança, Cachoeiro do Itapemirim, Cariacica, Castelo, Colatina, Conceição da Barra, Domingos Martins, Fundão, Iconha, João Neiva, Pancas, Pinheiros, Santa Maria, Santa Teresa e Vila Velha.
A elaboração do \textbf{Currículo} do Espírito Santo fundamenta -se em documentos legais que legitimam as \textbf{políticas} públicas educacionais, como : - Declaração Universal dos Direitos Humanos, \textbf{publicada} 1948, cujo documento o Brasil é signatário, assumindo o compromisso internacional pela educação, em seu \textbf{artigo} 26 estabelece que : > A instrução será orientada no sentido do pleno desenvolvimento da personalidade humana e do \textbf{fortalecimento} do respeito pelos direitos do ser humano e pelas liberdades fundamentais.
A instrução promoverá a compreensão, a tolerância e a amizade entre todas as nações e grupos raciais ou religiosos e coadjuvará as atividades das Nações Unidas em prol da manutenção da paz ( UNESCO, 1948 ). - Constituição Federal de 1988, em seu \textbf{Artigo} 205, determina : > A educação, direito de todos e dever do Estado e da família, será promovida e incentivada com a colaboração da sociedade, visando ao pleno desenvolvimento da pessoa, seu preparo para o exercício da cidadania e sua \textbf{qualificação} para o trabalho ( BRASIL, 1988 ). - Estatuto da Criança e do Adolescente ( Lei 8.069/1990 ), que dispõe sobre a proteção integral à criança e ao \textbf{adolescente}, definidos como pessoas em desenvolvimento físico, mental, moral, espiritual e social, que têm prioridade nas ações de proteção, de promoção e de defesa dos seus direitos, sem distinção de raça, cor ou classe social, e acrescenta em seu \textbf{Artigo} 4º : > É dever da família, da comunidade, da sociedade em geral e do poder público assegurar, com absoluta prioridade, a efetivação dos direitos referentes à vida, à saúde, à alimentação, à educação, ao esporte, ao lazer, à profissionalização, à cultura, à dignidade, ao respeito, à liberdade e à convivência familiar e comunitária ( BRASIL, 1990 ). - Estatuto da Juventude ( Lei 12.852/2013 ), que dispõe sobre os direitos dos jovens de 15 a 29 anos e declara, em seu \textbf{Artigo} 7º, a necessidade de \textbf{garantia} de educação básica, obrigatória e gratuita inclusive para os que a ela não tiveram acesso na idade adequada e complementa : > § 2º É dever do Estado oferecer aos jovens que não concluíram a educação básica programas na modalidade da educação de jovens e adultos, adaptados às necessidades e especificidades da juventude, inclusive no período noturno, ressalvada a \textbf{legislação} educacional específica ( BRASIL, 2013 ). - Lei de \textbf{Diretrizes} e Bases da Educação Nacional ( Lei 9394/96 ), em seu inciso IV, Art. 9º, afirma que cabe à União : > estabelecer, em colaboração com os Estados, o Distrito Federal e \textbf{os} Municípios, competências e \textbf{diretrizes} para a Educação Infantil, o Ensino Fundamental e o Ensino Médio, que nortearão os \textbf{currículos} e seus conteúdos mínimos, de modo a assegurar formação básica comum ( BRASIL, 1996 ). - Parâmetros Curriculares Nacionais, \textbf{publicados} em 1997, especificam que : > [... ] na medida em que o princípio da equidade reconhece a diferença e a necessidade de haver condições diferenciadas para o processo educacional, tendo em vista a \textbf{garantia} de uma formação de qualidade para todos, o que se apresenta é a necessidade de um referencial comum para a formação escolar no Brasil, capaz de indicar aquilo que deve ser garantido a todos, numa realidade com características tão diferenciadas, sem promover uma uniformização que descaracterize e desvalorize \textbf{peculiaridades} culturais e regionais ( MEC/SEF, 1997, p. 28 ). - \textbf{Diretrizes} Curriculares Nacionais Gerais da Educação Básica, Resolução CNE/CEB Nº 4/2010, que estabelecem em seu \textbf{Artigo} 13, § 3º : > A organização do percurso formativo, aberto e contextualizado, deve ser construída em função das \textbf{peculiaridades} do meio e das características, interesses e necessidades dos estudantes, incluindo não só os componentes curriculares centrais obrigatórios, \textbf{previstos} na \textbf{legislação} e nas normas educacionais, mas outros, também, de modo \textbf{flexível} e variável, conforme cada projeto escolar [... ] ( BRASIL, 2010 ). - \textbf{Diretrizes} Curriculares Nacionais para a Educação Infantil, Resolução CNE/CEB Nº 5/2009, que em seu \textbf{Artigo} 3º conceituam o \textbf{currículo} como : > [... ] conjunto de práticas que buscam articular as experiências e os saberes das crianças com os conhecimentos que fazem parte do patrimônio cultural, artístico, ambiental, científico e tecnológico, de modo a promover o desenvolvimento integral de crianças de 0 a 5 anos de idade ( BRASIL, 2009 ). - \textbf{Diretrizes} Curriculares Nacionais da Educação Básica para as modalidades da Educação do Campo ( Resolução CNE/CEB Nº 2/2008 ), da Educação Especial ( Resolução CNE/CEB Nº 4/2009 ), da Educação de Jovens e Adultos em contexto escolar ( Resolução CNE/CEB Nº 3/2010 ) e em privação de liberdade ( Resolução CNE/CEB Nº 2/2010 ), da Educação Escolar Indígena ( Resolução CNE/CEB Nº 5/2012 ), dos estudantes em situação de itinerância ( Resolução CNE/CEB Nº 3/2012 ), da Educação Escolar Quilombola ( Resolução CNE/CEB Nº 8/2012 ), que estabelecem as especificidades a serem \textbf{atendidas} em cada modalidade da educação básica nacional. - Resolução CEE/ES 3777/2014, em seu Art. 71, reconhece que : > O \textbf{currículo}, por ser uma construção social relacionada à ideologia, à cultura e à produção de identidades, tem ação direta na formação e no desenvolvimento dos estudantes, devendo, a sua elaboração privilegiar as seguintes relações : > I – cultura, sociedade e homem/mundo ; > II – conhecimento, produção de saberes e aprendizagem ; e > III – teoria e prática. - Plano Nacional de Educação, promulgado pela Lei nº 13.005/2014, \textbf{reitera} a necessidade de estabelecer e implantar, mediante pactuação interfederativa ( União, Estados, Distrito Federal e Municípios ), \textbf{diretrizes} pedagógicas para a educação básica e a base nacional comum dos \textbf{currículos}, com direitos e objetivos de aprendizagem e desenvolvimento dos(as) alunos(as) para cada ano do Ensino Fundamental e Médio, respeitadas as diversidades regional, estadual e local ( BRASIL, 2014 ). - A Base Nacional Comum Curricular ( BNCC ), \textbf{homologada} pela Resolução CNE/CP Nº 2, de 22 de dezembro de 2017 ( ) \textbf{Institui} e orienta a implantação da Base Nacional Comum Curricular, a ser respeitada \textbf{obrigatoriamente} ao longo das etapas e respectivas modalidades no âmbito da Educação Básica.
A BNCC trata das aprendizagens essenciais que todos os alunos devem desenvolver ao longo das etapas e modalidades da Educação Básica, de modo a que tenham assegurados seus direitos de aprendizagem e desenvolvimento, em \textbf{conformidade} com o que preceitua o Plano Nacional de Educação ( PNE ).
Este documento normativo aplica -se exclusivamente à educação escolar, tal como a define o § 1º do \textbf{Artigo} 1º da Lei de \textbf{Diretrizes} e Bases da Educação Nacional ( LDB, Lei nº 9.394/1996 ), e está orientado pelos princípios éticos, políticos e estéticos que visam à formação humana integral e à construção de uma sociedade justa, democrática e inclusiva, como fundamentado nas \textbf{Diretrizes} Curriculares Nacionais da Educação Básica ( DCN ). - Lei complementar Nº 799, de 12 de junho de 2015, cria o Programa de Escolas Estaduais de Ensino Médio em Turno Único, com o objetivo de planejar, executar e avaliar um conjunto de ações inovadoras em conteúdo, método e \textbf{gestão}, direcionadas à melhoria da \textbf{oferta} e da qualidade do ensino médio na rede pública do Estado, assegurando a criação e a implementação de uma rede de Escolas de Ensino Médio em Turno Único. - Pacto de Aprendizagem do Espírito Santo, Lei Nº 10.631, de 28 de março de 2017, que tem por objetivo \textbf{viabilizar} e fomentar o regime de colaboração entre a rede estadual e as redes \textbf{municipais} de ensino, a partir do diálogo permanente e ações conjuntas voltadas ao \textbf{fortalecimento} da aprendizagem e à melhoria dos indicadores educacionais dos alunos, das unidades de ensino e das referidas redes da educação básica no Espírito Santo, envolvendo domínio de competências de leitura, escrita e cálculo, adequados a cada idade e escolarização nas duas primeiras etapas de ensino da educação básica.
Os documentos supracitados respaldam a elaboração do \textbf{Currículo} do Espírito Santo, que tem como princípios o pleno desenvolvimento da pessoa, o exercício da cidadania, a \textbf{qualificação} para o trabalho, a equidade e a valorização das diferenças, a partir dos diversos contextos em que se configura a educação do nosso Estado.
A partir das aprendizagens essenciais definidas na BNCC, as habilidades foram contextualizadas, aprofundadas e complementadas considerando os sujeitos que estão implicados na educação do território do Espírito Santo.
Para sua concretização, foi essencial o regime de colaboração entre Estado e \textbf{municípios}, e demais parceiros.
Isso equivale a compreender o \textbf{currículo} como construção histórica e social.

% matched lemmas: adolescente, artigo, atender, conformidade, currículo, diretriz, equipe, flexível, fortalecimento, garantia, gestão, homologar, instituir, legislação, municipal, município, obrigatoriamente, oferta, os, peculiaridade, política, prever, publicar, qualificação, reiterar, viabilizar
\end{textsample}
